
\paragraph{\textit{Proof.}}
We may suppose that \(\mathfrak g\) has constant rank \(n\) and
constant nilpotent rank \(r\). The assertions concerning \(X\) in
\textup{a)} and \textup{b)} follow immediately from those concerning
\(\mathcal D\) and from the fact that \(X\) is a vector bundle over
\(\mathcal D\), defined by a locally free module; they were included
only for completeness.

\medskip
\noindent\textup{a)}
The functor \(\mathcal D\) is a subfunctor of
\(\operatorname{Grass}_{n-r}(\mathfrak g)\), whose value at \(S'\)
is the set of locally free quotient modules of rank \(n-r\) of
\(\mathfrak g_{S'}\). It is well known that the latter functor is
represented by a projective prescheme smooth over \(S\); see, for
example, Cartan Seminar 1960/61, Exposé~12, Nos.~2 and~3, whose
constructions carry over unchanged to preschemes.\begingroup
\renewcommand{\thefootnote}{*}\footnote{See also EGA~I, second
edition \textup{(announced as forthcoming from North-Holland
Publishing Company)}.}
\endgroup\footnote{Editors' note: See \S \SGARefLink{sga3:xiv:1.3}{I.1.3} of M.~Demazure and
P.~Gabriel, \emph{Groupes algébriques}, Masson \textup{(1970)}.}

We are therefore reduced to the following relative problem. Given
a locally free quotient module of rank \(n-r\) of \(\mathfrak g\),
or, equivalently, a locally free submodule \(\mathfrak d\) of rank
\(r\) that is locally a direct factor, represent the functor
\[
 F(S')=
 \begin{cases}
   \varnothing,
     &\text{if \(\mathfrak d_{S'}\) is not a Cartan subalgebra of
       \(\mathfrak g_{S'}\),}\\
   \{\varnothing\},
     &\text{otherwise.}
 \end{cases}
\]
We shall show that \(F\) is represented by a subprescheme of finite
presentation over \(S\). This will show that
\[
  \mathcal D\longrightarrow
  \operatorname{Grass}_{n-r}(\mathfrak g)
\]
is represented by an immersion of finite presentation and will
complete the proof of \textup{a)}.

First express the condition that \(\mathfrak d_{S'}\) be a Lie
subalgebra of \(\mathfrak g_{S'}\). This condition says precisely
that \(S'\to S\) factors through a certain closed subprescheme
\(S_1\) of finite presentation over \(S\); locally on \(S\), its
equations are obtained at once from a basis of \(\mathfrak g\)
adapted to \(\mathfrak d\). We may therefore suppose that
\(S=S_1\).

Next we must express that \(\mathfrak d_{S'}\) contains,
fpqc-locally, a quasi-regular section of \(\mathfrak g_{S'}\). For
this purpose consider
\[
  V=W(\mathfrak d)\cap U,
\]
where \(U\) is the open subset of regular points of
\(W(\mathfrak g)\) \textup{(\SGARefLink{sga3:xiv:2.10}{2.10})}. Since the structure morphism
\(V\to S\) is smooth and quasi-compact, its image \(S_1\) is an
open subset of \(S\), and the immersion \(S_1\to S\) is
quasi-compact, hence of finite presentation. The condition on
\(S'\) is exactly that \(S'\to S\) factor through \(S_1\). We are
thus reduced to the case \(S=S_1\), and, by descent, to the case in
which \(\mathfrak d\) has a section \(a\) that is quasi-regular in
\(\mathfrak g\).

It remains to express that, for the induced section \(a_{S'}\), the
endomorphism
\[
  \operatorname{ad}(a_{S'})
    _{\mathfrak g_{S'}/\mathfrak d_{S'}}
\]
is bijective. This again says that \(S'\to S\) factors through an
open subprescheme of finite presentation, namely \(S_D\), where
\[
  D=\det\!\left(
    \operatorname{ad}(a)_{\mathfrak g/\mathfrak d}
  \right).
\]
But \(\mathfrak d|_{S_D}\) is then a Cartan subalgebra of
\(\mathfrak g|_{S_D}\). Consequently \(S_D\) represents \(F\),
which proves \textup{a)}.

\medskip
\noindent\textup{b)}
This follows immediately from \SGARefLink{sga3:xiv:2.11}{2.11}\textup{(a)} and \SGARefLink{sga3:xi:remark:1:5}{XI~1.5}. It is
also, of course, a consequence of the more precise assertion
\textup{c)}.

\medskip
\noindent\textup{c)}
Let \(a_0\) be a regular point of \(\mathfrak g(s)\) contained in
\(\mathfrak d_0\), and extend it to a section \(a\) of
\(\mathfrak g\) over an open neighborhood of \(s\). We may plainly
suppose that this neighborhood is all of \(S\). Choose a complement
\(M_0\) of \(\mathfrak d_0\) in the vector space
\(\mathfrak g(s)\). After again replacing \(S\) by an open
neighborhood of \(s\), there is a submodule \(M\) of
\(\mathfrak g\), locally a direct factor, such that
\(M(s)=M_0\); we again suppose that the neighborhood is \(S\).

Let \(\mathcal V\) be the subfunctor of \(\mathcal D\) for which
\(\mathcal V(S')\) consists of the Cartan subalgebras
\(\mathfrak d'\) of \(\mathfrak g_{S'}\) satisfying:
\begin{enumerate}
\item[\textup{1\textdegree)}]
  \(\mathfrak d'\) is a complement of \(M_{S'}\);
\item[\textup{2\textdegree)}]
  the unique section in
  \((a_{S'}+M_{S'})\cap\mathfrak d'\) is a regular section of
  \(\mathfrak g_{S'}\).
\end{enumerate}
Condition \textup{1\textdegree)} defines an open \(V_1\) of
\(\mathcal D\), induced by the open of
\(\operatorname{Grass}_{n-r}(\mathfrak g)\) given by the same
condition. By \SGARefLink{sga3:xiv:2.10}{2.10} and \((\mathrm C_2)\), the conjunction of
\textup{1\textdegree)} and \textup{2\textdegree)} defines an open
of \(V_1\). Thus \(\mathcal V\) is represented by an open
subprescheme \(V\) of \(\mathcal D\), plainly containing the point
\(d\).

Now let \(\mathcal V'\) be the subfunctor of \(W(M)\) defined by
\[
 \mathcal V'(S')
 =
 \left\{
   u'\in\Gamma(S',M_{S'})
   \ \middle|\
   \begin{array}{l}
   \textup{(i) }a_{S'}+u'\text{ is a regular section of }
     \mathfrak g_{S'},\\
   \textup{(ii) the unique Cartan subalgebra }\mathfrak d'
     \text{ of }\mathfrak g_{S'}\\
   \quad\text{containing }a_{S'}+u'
     \text{ is a complement of }M_{S'}
   \end{array}
 \right\}.
\]
Condition \textup{(i)} defines an open subprescheme \(V'_1\) of
\(W(M)\): it is the inverse image of the open \(U\) of regular
points of \(W(\mathfrak g)\) \textup{(\SGARefLink{sga3:xiv:2.10}{2.10})} under the translation
morphism \(m\mapsto a+m\). The conjunction of \textup{(i)} and
\textup{(ii)} defines an open \(V'\) of \(V'_1\), namely the
inverse image of \(V\) under the evident morphism
\[
  V'_1\longrightarrow\mathcal D
\]
that associates with \(u'\) the unique Cartan subalgebra
\(\mathfrak d'\) containing \(a_{S'}+u'\). Its restriction gives
a morphism
\[
  V'\longrightarrow V,
\]
which is plainly an isomorphism. This proves \textup{c)}.

\paragraph{\textbf{Corollary 2.17.}}
\SGARefTarget{sga3:xiv:2.17}
{\itshape Let \(\mathfrak g\) be a finite-dimensional Lie algebra
over a field \(k\). Then the scheme \(\mathcal D\) of Cartan
subalgebras of \(\mathfrak g\) \textup{(\SGARefLink{sga3:xiv:2.16}{2.16}\textup{(a)})} is
quasi-projective, smooth, and irreducible. If \(\mathfrak g\)
contains a regular element \textup{(for example, if \(k\) is
infinite)}, then \(\mathcal D\) is a rational variety; that is, its
function field is a purely transcendental extension of \(k\).}

Indeed, \(\mathcal D\) is irreducible because there is a surjective
morphism
\[
  \psi^{-1}(U)\longrightarrow\mathcal D,
\]
and \(\psi^{-1}(U)\) is irreducible, being isomorphic to the open
subset \(U\) of \(W(\mathfrak g)\). The assertion concerning the
function field follows immediately from \SGARefLink{sga3:xiv:2.16}{2.16}\textup{(c)}.

\paragraph{\textbf{Remarks 2.18.}}
\SGARefTarget{sga3:xiv:2.18}
I do not know whether this conclusion remains true when \(k\) is
finite without assuming that \(\mathfrak g\) contains a regular
point; compare \SGARefLink{sga3:xiv:2.12}{2.12}. It can be proved when \(\mathfrak g\) is the
Lie algebra of an algebraic group \(G\) smooth over \(k\), at least
when \(G/\operatorname{radical}(G)\) is an ``adjoint'' semisimple
group, by using a result of Chevalley cited below \textup{(see the
Appendix)}. It is plausible that the result remains valid without
restriction on \(G\); for this it would suffice that the cited
result of Chevalley be proved for every semisimple algebraic group,
not necessarily an adjoint group.
